\section{Discussion and Limitations}\label{sec:discussion}

\textbf{What is real and what is modelled.} The benign and attack traffic are real N-BaIoT
captures~\cite{nbaiot}, and the C implementation is real and compiled for the ESP32. Three things are
ours and synthetic: the concurrency of the fleet (the nine devices were captured separately), the honest
drift (physically modelled transforms of the Kitsune statistics, not a recorded network change), and the
poisoner (crafted feature vectors, not packets). A public intrusion dataset with documented concept
drift did not exist when the 2025 survey of intrusion datasets was written~\cite{nidsurvey}, which is
why the drift is injected; the testbed's real common-mode drift is the next step.

\textbf{The frontier cannot be escaped, only placed.} Proposition~\ref{prop:frontier} is a limit on
every rule, not a weakness of FOA: a device's own drift and an attack on that device look the same to
everyone else. FOA's contribution is to move all drift the peers can vouch for out of that dilemma and
to make the remaining trade-off one explicit number, $B$. A deployment that expects frequent device-only
changes (firmware updates) can raise $B$ for a bounded window when the update is signed and announced, as
licensed-update schemes do; we leave that to future work.

\textbf{When the fleet is the attacker.} If one exploit compromises a majority of a device's peer group,
or a gateway rewrites everyone's traffic, consensus carries the attack. Vulnerability-diverse peers make
the first harder, not impossible: N-BaIoT itself shows that one credential-guessing botnet (BASHLITE)
infected all nine devices~\cite{nbaiot}. FOA then degrades to its residual budget and common-channel
budget, which still bound the rate of movement (Theorem~\ref{thm:byz}), but not to safety.

\textbf{Limits we measured.} Three results are against FOA and we report them as such. Byzantine
peers cannot compromise a device on their own, but together with a traffic poisoner they help it: with
three of seven peers controlled, about half of the victims fell during strong drift, all of them victims
whose attack traffic sits close to their benign traffic. A device with a large drift of its own on
top of the shared one is lost by FOA (and by every budgeted variant), while single-device adaptation
keeps it. Core-hiding and mimicry attackers at share 0.4 still compromised about a third of victims,
mostly the two cameras whose scan traffic is closest to benign.

\textbf{Spread is the attacker's lever.} In the compromised runs we traced, the successful attacks
inflated the spread rather than dragging the centre, the radius attack Kloft and Laskov warned
about~\cite{kloft}. FOA's spread budget is what holds it, and it is the parameter we would tune first.

\textbf{Hardware.} The boards had not arrived when this paper was written. The firmware is complete and
checked against the reference; latency, energy and radio loss on real ESP32s, and the fleet basis fitted
on the testbed's own traffic, are reported as host measurements and estimates here and will be replaced
by measurements.
